% ============================================================================
%  Subdocumento 1. Contenido. Se usa igual en el documento único y en el
%  PDF propio del subdocumento. Los formularios que le asigna el T-21 se
%  agregan solos al final (datos en formularios/ de esta carpeta).
% ============================================================================

\begin{resumenApertura}
\marcador{Qué resuelve este subdocumento, en cuatro o cinco líneas}
\begin{recibe}
  \item \marcador{Qué recibe el mandante}
\end{recibe}
\end{resumenApertura}

\section{Trayectoria, capacidades instaladas y servicios}\label{sec:1-trayectoria-capacidades-instalad}

\marcador{Reseña de la trayectoria, capacidades instaladas, líneas de negocio, productos y servicios ofrecidos}

\section{Estructura organizacional, dotación, certificaciones y alianzas}\label{sec:1-estructura-organizacional-dotaci}

\marcador{Estructura, dotación, certificaciones institucionales y alianzas tecnológicas vigentes}

\section{Experiencia en la industria y en proyectos equivalentes}\label{sec:1-experiencia-en-la-industria-y-en}

\marcador{Experiencia en transporte de carga y en proyectos de complejidad equivalente, con volúmenes comparables. El detalle va en el Formulario T-6}

\section{Modelo de gobierno de calidad, seguridad y conocimiento}\label{sec:1-modelo-de-gobierno-de-calidad-se}

\marcador{Roles, comités, indicadores y ciclo de auditoría interna}

\section{Credenciales financieras}\label{sec:1-credenciales-financieras}

\marcador{Antecedentes que sostienen 56 meses de contrato}
